\section{Polygon reconstruction, winding, and triangle emission}
\label{sec:polygon}

\subsection{From four quadrants to a ring}

For a sub-segment the shader resolves four quadrant cells.  Not all of
them are surface cells: a coarse cell can have no zero crossing at the
sample.  The algorithm:
\begin{enumerate}
\item iterate quadrants in the fixed order $q_0,\ldots,q_3$;
\item skip non-surface cells;
\item compute the cell's dual vertex position with the memoized QEF;
\item reject a cell if its node id or position duplicates an already
  collected cell (within the edge epsilon);
\item stop at four collected cells.
\end{enumerate}
The deduplication is not merely an optimization.  A coarse ancestor can
appear in two opposite quadrants of the ring, especially on world-axis
cell boundaries where the sample nudge stays inside the same large cell.
The old code removed only adjacent duplicates and then aborted the whole
segment if any duplicate remained; that is exactly why triangles were
missing along the $x=0$, $y=0$, and $z=0$ planes.  Deduplicating against
all collected cells keeps the remaining 3 or 4 corners and preserves their
order.

\subsection{Rotating the owner to the front}

The ring is rotated so the owner cell is first:
\begin{lstlisting}[style=glsl,caption={Owner rotation.},label={lst:rotate}]
find k with ring[k].id == owner.id or ring[k].pos == owner.pos;
if (k > 0) rotate ring left by k;
\end{lstlisting}
Both chunks on either side of a boundary resolve the same owner (the
ordering key is global), so their duplicated boundary quads get the same
rotation.  Without this, the same quad would be triangulated from different
anchors on the two sides, producing different triplanar planes and
z-fighting along the seam.

\subsection{Three-ring and four-ring fans}

A three-ring emits one triangle:
\begin{equation}
  (r_0,r_1,r_2).
\end{equation}
A four-ring emits two:
\begin{equation}
  (r_0,r_1,r_2), \quad (r_0,r_2,r_3).
\end{equation}
If the owner's start sign is positive (outside), the order is reversed:
\begin{equation}
  (r_0,r_2,r_1), \quad (r_0,r_3,r_2).
\end{equation}
This is a fan around the owner.  A ring with fewer than three cells emits
nothing; it is a boundary case where the surface only touches one or two
quadrants.

\subsection{The winding oracle}

The winding is not derived from the triangle normal.  It is derived from
the SDF sign at the segment start, evaluated with the owner's trilinear
field:
\begin{equation}
  \text{solidAtStart} = \bigl(d_{\text{owner}}(a) < 0\bigr).
\end{equation}
If the owner's interpolation returns $+\infty$, the producing cell's
endpoint sign is used.  The producing cell is the cell that scanned the
edge, and it only scans edges whose own corner signs differ; those signs
are therefore authoritative for existence.

Why not use the QEF normal?  At coarse LoD a cell can have only one or two
crossing edges, and the QEF normal is poorly constrained.  A triangle
wound from a noisy normal can flip at a LoD boundary, making the surface
black or invisible from one side.  The sign oracle is independent of the
normal and robust across cell sizes.

\subsection{Degenerate triangles}

Before appending, the shader compares the three dual positions:
\begin{equation}
  \|p_i-p_j\| \le \epsilon \quad\Rightarrow\quad \text{drop}.
\end{equation}
This catches two cases:
\begin{itemize}
\item two cells collapsed to the same QEF position (a thin sliver);
\item a cell with an exact-zero corner snapped to the corner, coinciding
  with a neighbouring cell's vertex.
\end{itemize}
A zero-area triangle would break normals, shadow-map rasterization, and
ray-triangle intersections.  Dropping it is safer than emitting it.

\subsection{Direct index append and lazy vertices}

The triangle is appended straight into the chunk's index span:
\begin{lstlisting}[style=glsl,caption={Index append.},label={lst:indexappend}]
uint base = atomicAdd(indexCount, 3u);
int s0 = nodeSlot(c0.id), s1 = nodeSlot(c1.id), s2 = nodeSlot(c2.id);
if (s0 < 0 || s1 < 0 || s2 < 0) { atomicAdd(invalidSlotCount, 1u); return; }
uint extent = max(max(s0, s1), s2) + 1u;
atomicMax(vertexExtent, extent);
if (base + 3u > pc.indexCapacity || extent > pc.vertexCapacity) {
    atomicAdd(indexOverflow, (base + 3u > pc.indexCapacity) ? 1u : 0u);
    atomicAdd(vertexOverflow, (extent > pc.vertexCapacity) ? 1u : 0u);
    atomicAdd(droppedCount, 1u);
    return;
}
poolIdx[pc.indexBase + base + 0u] = uint(s0);
poolIdx[pc.indexBase + base + 1u] = uint(s1);
poolIdx[pc.indexBase + base + 2u] = uint(s2);
writeVertexOnce(c0.id, s0);
writeVertexOnce(c1.id, s1);
writeVertexOnce(c2.id, s2);
\end{lstlisting}
The atomic only reserves the base; the three index writes are private.
The counter advances by three even when the triangle does not fit, so the
host reads the exact index requirement for the overflow-growth retry.  No
write below the capacity check is ever outside the reserved span, so an
overflow cannot corrupt another chunk.

The vertex payload is a pure function of its node.  Many corners reference
the same surface node, so the shader elects one writer per node with an
\code{atomicExchange} on the device-local \code{nodeWritten} buffer, which
is cleared before every dispatch.  The elected invocation writes the
64-byte payload and increments \code{vertexWriteCount}; the others skip it.
By the time the dispatch completes, every referenced payload is therefore
written exactly once.

\subsection{Triangle order and visual invariance}

The order of triangles in the index stream is the order in which scan cells
happen to append.  Different GPUs, workgroup schedules, or occupancy can
produce a different order.  This is acceptable because:
\begin{itemize}
\item the triangle set is identical (each cell's output is independent);
\item the renderer draws the whole index range, not a prefix;
\item the shadow map, depth prepass, and ray tracing are order-independent;
\item the vegetation and fire anchors are keyed by chunk/node, not by
  triangle index.
\end{itemize}
The only place order matters is debugging: a diff of two index streams
must be order-insensitive.  The engine's self-check compares triangle sets,
not raw arrays.

\subsection{The boundary quad}

A duplicated boundary quad is emitted by both neighbouring chunks.  Both
copies must be identical in position, order, and material.  The engine
guarantees this by:
\begin{enumerate}
\item using the same owner ordering key (size, min, node index);
\item sampling the same four quadrants (the neighbour's node table contains
  the same cells);
\item using the same QEF arithmetic and the same edge epsilon;
\item rotating the owner to the front.
\end{enumerate}
The duplicate triangles are not removed at the chunk level; the renderer
draws both.  They coincide exactly, so the depth test resolves the
z-fight.  Removing them would require a cross-chunk dedup pass and would
reintroduce a seam if the dedup key differed.
